% Full proofs for function-information and differentiability constructions.

\begin{annexitem}{annex:p04-ca-sample-modulus}{thm:p04-ca-sample-modulus}{Both directions between uniform and sample-and-modulus names}
\begin{proof}
For a uniform polygon name \((P_j)\), the rational value
\(P_{n+1}(q_i)\) has error at most \(2^{-n-1}\) at \(q_i\).
Each finite rational polygon is evaluated exactly by comparing rational
arguments with its rational vertices. Enumerating the pairs \((i,n)\)
thus gives a productive sample table. The modulus algorithm of
\cref{thm:p04-fo-evaluation} reads \(P_{n+2}\), computes its rational
maximum absolute slope \(L\), and selects an integer \(\mu(n)\)
with \(L2^{-\mu(n)}\leq2^{-n-1}\). The two polygon-error terms
then give the required modulus bound. This computes both components,
not merely the existence of the second.

Conversely, fix precision \(n\), read \(\mu(n+2)\), and put
\(N=2^{\mu(n+2)}\). The partition vertices are \(t_j=j/N\),
\(0\leq j\leq N\). Search the effective enumeration \((q_i)\)
until each of these finitely many rational numbers has occurred;
rational equality is decidable, and surjectivity ensures termination.
Read the corresponding table approximations at accuracy \(n+2\)
and call them \(y_j\). Produce the rational polygon with vertices
\((t_j,y_j)\).
For \(x\) in a partition cell \([t_j,t_{j+1}]\), write
\(x=(1-a)t_j+at_{j+1}\), where \(0\leq a\leq1\).
Both endpoints are within \(2^{-\mu(n+2)}\) of \(x\), hence
\[
 \begin{split}
 |(1-a)y_j+ay_{j+1}-f(x)|
 &\leq (1-a)|y_j-f(t_j)|+a|y_{j+1}-f(t_{j+1})|\\
 &\quad +(1-a)|f(t_j)-f(x)|+a|f(t_{j+1})-f(x)|\\
 &\leq2^{-n-2}+2^{-n-2}=2^{-n-1}.
 \end{split}
\]
Every output stage uses finitely many entries and terminates for every
valid sample-and-modulus name. The resulting polygons therefore meet
the actual \(\gamma_{\mathrm P}\) accuracy clause. The two constructions
decode to the same original function, proving computable equivalence.
\end{proof}
\end{annexitem}

\begin{annexitem}{annex:p04-ca-sample-obstruction}{thm:p04-ca-sample-obstruction}{A missed tent with a fixed area}
\begin{proof}
Let \(p\) be the sample name of the zero function that uses a fixed
rational code for zero in every table cell. Suppose a continuous area
realizer exists. Its scalar output at \(p\) names zero, so its
approximation \(u_3\) at precision three satisfies \(|u_3|\leq1/8\).
Continuity fixes that entry on some finite prefix of \(p\).
Only finitely many rational sample locations occur in those table
positions. Choose rational \(0<a<b<1\) with \([a,b]\)
disjoint from that finite set, and set \(c=(a+b)/2\).
The rational polygon \(T\) with vertices
\[
 (0,0),\ (a,0),\ (c,2/(b-a)),\ (b,0),\ (1,0)
\]
is zero at all observed locations and has integral one: the triangle's
base has length \(b-a\), its height is \(2/(b-a)\), and its area
is half their product.
It has a valid sample name with the identical finite zero prefix.
Indeed, all remaining table entries can be its exact rational values
at rational arguments. The area realizer must use the same \(u_3\)
there, although now \(|u_3-1|\leq1/8\) is required. The two
inequalities are incompatible.

If a continuous translator to polygon names existed, composing it
with the computable integration realizer of
\cref{thm:p04-fo-integration} would yield the just-forbidden continuous
area realizer. Finite-use continuity similarly rules out computable
translators. This argument works with exact rational answers at every
queried location; approximation noise is not its cause.
\end{proof}
\end{annexitem}

\begin{annexitem}{annex:p04-ca-computable-c1}{thm:p04-ca-computable-c1}{Disjoint supports, continuity at zero and a uniform polygon algorithm}
\begin{proof}
On \([-1,1]\), differentiation of \(\theta\) gives
\(\theta'(x)=5x^4-6x^2+1\). At the endpoints both the value and
this derivative are zero, matching the zero pieces outside, so
\(\theta\in C^1(\mathbb R)\). Setting \(r=x^2\in[0,1]\),
the derivative polynomial \(5r^2-6r+1\) has minimum \(-4/5\)
at \(r=3/5\) and maximum one; hence \(|\theta'|\leq1\).
Its zeros in the interior are \(\pm1/\sqrt5\), where
\(\theta(1/\sqrt5)=16/(25\sqrt5)\); oddness and the sign on
each half of the support prove the stated extrema and amplitude.
The chain rule proves all scaled-bump clauses.

The positive constant \(\sqrt5\) has a computable scalar name:
bisect \([2,3]\) using rational comparisons of the midpoint's square
with five. The nested intervals contain \(\sqrt5\) and have known
geometrically decreasing lengths. Their midpoints give the required
Cauchy approximations. Scalar arithmetic and the nonzero reciprocal
algorithm then compute \(\lambda\). Moreover
\(\sqrt5<9/4\) implies \(\lambda>64/225>1/4\), so
\(1/\lambda<4\).

The support with centre \(2^{-n}\) has radius
\(\alpha_n\leq2^{-2n-2}\leq2^{-n-2}\).
For consecutive possible centres \(2^{-n}\) and \(2^{-n-1}\),
the sum of these radius bounds is at most
\(2^{-n-2}+2^{-n-3}=3\cdot2^{-n-3}\), smaller than their
distance \(2^{-n-1}\). Thus their closed support intervals
are disjoint, and all other pairs are disjoint as well: each interval
lies between its two neighbouring separating gaps. Every compact
subinterval of \((0,1]\) meets only finitely many supports.
The function is therefore locally a finite \(C^1\) sum away from zero,
including the matching zero derivatives at support boundaries.

If \(x>0\) belongs to support \(n\), then
\(x\geq2^{-n}-\alpha_n\geq(3/4)2^{-n}\), and
\[
 \frac{|f(x)|}{x}
 \leq \frac{\beta_n}{(3/4)2^{-n}}
 =\frac43\alpha_n.
\]
As \(x\to0\) through supports, their indices tend to infinity,
so this tends to zero. Outside them the quotient is zero.
Since \(f(0)=0\), the right derivative exists and is zero.
Within support \(n\), the derivative's magnitude is at most
\(\beta_n/(\lambda\alpha_n)=2^{-n}/\lambda\), which also
tends to zero. This proves continuity of \(f'\) at zero, and hence
\(f\in C^1[0,1]\).
Each point lies in at most one support, so \(|f'(x)|<4\).
The mean value theorem gives the Lipschitz bound four.
At the centre, the bump's derivative is exactly \(2^{-n}/\lambda\)
when present. If it is absent, no other support contains that centre,
and the local derivative is zero. This proves the centre formula.

An undecidable subset of \(\mathbb N\) and its complement are both
infinite: a finite or cofinite set is decidable by a program with its
finite exception list. Along centres with \(n\in A\),
\(f'(2^{-n})/2^{-n}=1/\lambda\); along centres with \(n\notin A\)
the same quotient is zero. Both subsequences approach zero, and
\(f'(0)=0\), so the second right derivative at zero cannot exist.

After the first \(M+1\) enumeration terms, at any point the difference
\(f-g_M\) has either no nonzero term or just one with index \(k>M\).
Its magnitude is bounded by
\[
 \beta_{h(k)}=2^{-k-3h(k)-2}\leq2^{-M-3}.
\]
This is a computable uniform tail bound even though there is no
computable decision procedure for membership in \(A\).
All finite partial sums have the same derivative bound four.

Here is a polygon-name algorithm, not merely a pointwise evaluation
claim. At output precision \(n\), compute \(h(0),\ldots,h(n+1)\)
and form the finite sum \(g_{n+1}\).
Use the dyadic mesh \(d=2^{-n-4}\), and at every vertex compute a
rational approximation to \(g_{n+1}\) with error at most
\(2^{-n-3}\). This terminates by the clamping formula, finite scalar
arithmetic, and the computed positive constant \(\lambda\).
Interpolate these values linearly. On each cell, interpolation of
the exact endpoint values differs from \(g_{n+1}\) by at most
\(4d=2^{-n-2}\), by its Lipschitz bound and convex endpoint weights.
The approximate endpoint values add at most \(2^{-n-3}\).
Finally the uniform tail adds at most \(2^{-n-4}\).
The total is \((1/4+1/8+1/16)2^{-n}<2^{-n}\).
Thus the produced finite rational polygon meets the required name
clause at every output position. The algorithm only waits for finitely
many values of the total enumeration \(h\); it never tests a negative
membership assertion about \(A\).
\end{proof}
\end{annexitem}

\begin{annexitem}{annex:p04-ca-noncomputable-derivative}{thm:p04-ca-noncomputable-derivative}{A decidable gap would decide the enumerable set}
\begin{proof}
Suppose \(f'\) had a computable uniform polygon name.
Evaluation at the computable rational input \(2^{-n}\) would then
compute \(f'(2^{-n})\) uniformly in \(n\).
Because \(\lambda<1\), a present bump has derivative
\(2^{-n}/\lambda>2^{-n}\), whereas an absent bump gives zero.
Compute a rational approximation \(v\) to the derivative there with
error at most \(2^{-n-3}\).
In the absent case \(v\leq2^{-n-3}<2^{-n-1}\).
In the present case
\(v>2^{-n}-2^{-n-3}>2^{-n-1}\).
The rational comparison with \(2^{-n-1}\) therefore decides
\(n\in A\), a contradiction. The approximation never lies on
the decision threshold, so no real-equality test has been used.

If \(f'\) had a computable modulus, the already computable name of
\(f\), together with that modulus, would compute its derivative by
\cref{thm:p04-fo-derivative-promise}. This contradicts the first part.
Thus neither a derivative name nor a derivative modulus is computable.
For any one centre, its actual value is one of two computable constants;
that individual existence statement supplies no uniform selector for
the alternatives and does not refute the argument.
\end{proof}
\end{annexitem}

\begin{annexitem}{annex:p04-ca-derivative-modulus}{thm:p04-ca-derivative-modulus}{Both modulus directions, a fixed bound and higher orders}
\begin{proof}
If \(f'\) has a computable polygon name, applying the computed-modulus
construction of \cref{thm:p04-fo-evaluation} to that name gives a
computable derivative modulus. Conversely, if such a modulus is
computable, supply its stream together with the computable name of
\(f\) to \cref{thm:p04-fo-derivative-promise}.
That theorem's secant construction gives a uniform derivative name.
These establish both directions for the fixed computable function.

If \(f''\) is continuous, compactness gives a finite bound on its
absolute value. Choose a positive integer \(K\) at least that bound,
and an integer \(\ell\) with \(K\leq2^\ell\).
For \(x,y\in[0,1]\), the mean value theorem applied to \(f'\)
gives
\[
 |f'(x)-f'(y)|\leq K|x-y|.
\]
Consequently \(|x-y|\leq2^{-n-\ell}\) implies
\(|f'(x)-f'(y)|\leq2^{-n}\).
The integer function \(n\mapsto n+\ell\) is a computable modulus.
For one fixed \(f\), the existing integer \(\ell\) can be part of
the program; a claim to compute it from the ordinary input name has
not been made. If \(K\) is actually supplied, the rational search for
\(\ell\) terminates and gives the uniform stronger-input algorithm.

For a fixed computable \(C^\infty\) function, start with its own
computable name. The fixed \(C^2\) result gives a computable name of
its first derivative. That derivative is itself \(C^2\), so the
result applies again. Induction proves, for every fixed finite order,
existence of a computable name of that derivative. The proof makes a
new finite bound choice at each order and supplies no effective
sequence of those choices. It therefore does not establish uniform
computability of the derivative sequence.
\end{proof}
\end{annexitem}

\begin{annexitem}{annex:p04-ca-root-certificate}{thm:p04-ca-root-certificate}{The lower-slope certificate and its sensitivity constant}
\begin{proof}
The intermediate value theorem and the endpoint signs give a zero
\(z\). The mean value theorem gives, for \(x>y\),
\(f(x)-f(y)\geq m(x-y)>0\); thus \(f\) is strictly increasing
and its zero is unique. Applying the same inequality between \(x\)
and \(z\) gives \(|f(x)|\geq m|x-z|\), with the case \(x=z\)
included. This proves the residual-to-location estimate.

Enumerate rational \(s\in[0,1]\) and natural \(k\), evaluating
the finite rational polygon \(P_k\) at \(s\) exactly and testing
the displayed strict rational certificate. Each successful test gives
\[
 |f(s)|\leq|P_k(s)|+2^{-k}<m2^{-n},
\]
hence \(|s-z|<2^{-n}\).
To prove termination, continuity and rational density give a rational
\(s\in[0,1]\) with \(|f(s)|<m2^{-n}\), including endpoint
roots, where the root itself is rational. For that fixed \(s\), take
\(k\) sufficiently large that
\(2^{1-k}<m2^{-n}-|f(s)|\). Then
\[
 |P_k(s)|+2^{-k}
 \leq |f(s)|+2^{1-k}<m2^{-n}.
\]
The enumerated search therefore reaches a successful pair at every
output precision, producing a scalar root name on every promised input.

Finally, if \(g(w)=0\) and \(\lVert f-g\rVert_\infty\leq\epsilon\),
then \(|f(w)|=|f(w)-g(w)|\leq\epsilon\).
The already proved estimate gives \(|w-z|\leq\epsilon/m\).
No derivative or monotonicity promise on \(g\) is required for this
bound on each of its existing zeros; the assertion does not claim that
a perturbation necessarily retains a zero in the interval.
\end{proof}
\end{annexitem}
